% DERIVED from formal_layer.json — read-only, do not edit.
\begin{assumptionv}{ass:uniform-design}[Uniform covariate design]
The covariate \(X\) has the uniform probability law \(\lambda\) on the unit interval under the observed law \(P\):
\[
P\circ X^{-1}=\lambda.
\]
\end{assumptionv}

\begin{definitionv}{synth_2}[Log-odds map definition]
The log-odds map \(\operatorname{logit}\colon\mathbb R\to\mathbb R\) is defined by
\[
\operatorname{logit}(t)=\log\!\left(\frac{t}{1-t}\right),
\qquad t\in\mathbb R.
\]
Here division by zero has value zero, and the real logarithm is extended by \(\log 0=0\) and \(\log t=\log|t|\) for \(t<0\). For \(0<t<1\), this is the usual log-odds transformation.
\end{definitionv}

\begin{definitionv}{synth_4}[Native prognosis logit]
The native prognosis logit \(\nu(P,x)\) is the log odds of the conditional outcome risk under control, for \(x\in[0,1]\) and an observed probability law \(P\) with binary treatment and outcome whose conditional four-cell table satisfies:
\begin{itemize}
\item \textbf{(Conditional probabilities.)} At each covariate value, the four cells sum to one and give a conditional representation of \(P\) given the covariate.
\item \textbf{(Continuity and interiority.)} Each cell probability is continuous in the covariate and lies strictly between zero and one at every covariate value.
\item \textbf{(Full support.)} The covariate marginal assigns positive probability to every nonempty relatively open subset of \([0,1]\).
\end{itemize}
Its defining formula is
\[
\nu(P,x)
:=
\operatorname{logit}\!\left(
\frac{\Pr_P(\text{treatment}=0,\text{outcome}=1\mid X=x)}
{\Pr_P(\text{treatment}=0,\text{outcome}=0\mid X=x)
+\Pr_P(\text{treatment}=0,\text{outcome}=1\mid X=x)}
\right),
\qquad
\operatorname{logit}(t):=\log\!\left(\frac{t}{1-t}\right).
\]
\end{definitionv}

\begin{definitionv}{synth_1}[The logistic map]
The logistic map \(\ell:\mathbb R\to\mathbb R\) is defined by
\[
\ell(t)=\frac{1}{1+\exp(-t)},\qquad t\in\mathbb R.
\]
\end{definitionv}

\begin{assumptionv}{ass:homogeneous-logit}[Homogeneous conditional log-odds relation]
The treated-arm conditional outcome risk \(\mu_1(P,x)\) satisfies
\[
\mu_1(P,x)=\ell\{\nu(P,x)+\theta(P)\}
\qquad (x\in[0,1]),
\]
where \(\ell\) is the logistic map, \(\nu(P,x)\) is the native prognosis logit, and \(\theta(P)\) is the homogeneous conditional log-odds coefficient.
\end{assumptionv}

\begin{assumptionv}{ass:effect-envelope}[Compact scalar effect envelope]
For an observed law \(P\), the log-odds contrast at covariate zero,
\[
\theta(P)=\operatorname{logit}\bigl(\mu_1(P,0)\bigr)-\nu(P,0),
\]
satisfies
\[
|\theta(P)|\le \frac12.
\]
\end{assumptionv}

\begin{assumptionv}{ass:propensity-envelope}[Propensity logit envelope]
The native propensity-logit function \(g(P,\cdot)\) satisfies
\[
\|g(P,\cdot)\|_\infty\le 1.
\]
\end{assumptionv}

\begin{assumptionv}{ass:prognosis-envelope}[Prognosis logit envelope]
The native prognosis-logit function \(\nu(P,\cdot)\) satisfies
\[
\|\nu(P,\cdot)\|_\infty\le 1.
\]
\end{assumptionv}

\begin{assumptionv}{ass:propensity-holder}[Propensity Hölder seminorm bound]
The Hölder seminorm of the native propensity-logit function \(g(P,\cdot)\) at the public smoothness exponent \(\alpha\) satisfies
\[
[g(P,\cdot)]_\alpha\le 2.
\]
\end{assumptionv}

\begin{assumptionv}{ass:prognosis-holder}[Prognosis Hölder bound]
The native prognosis-logit function \(\nu(P,\cdot)\) satisfies the Hölder seminorm bound
\[
[\nu(P,\cdot)]_\beta\le 2.
\]
\end{assumptionv}

\begin{definitionv}{def:logistic-class}[Homogeneous logistic class \(\mathcal L\)]
\[
\mathcal L
=
\left\{
P:\mu_1(P,x)=\ell\{\nu(P,x)+\theta(P)\}
\ \text{for every }x\in[0,1]
\right\}.
\]
Here \(\mu_1(P,x)\) is the treated-arm conditional outcome risk, \(\ell\) is the logistic map, \(\nu(P,x)\) is the native prognosis logit, and \(\theta(P)\) is the homogeneous conditional log-odds coefficient. The defining relation is that of \cref{obj:ass:homogeneous-logit}.
\end{definitionv}

\begin{definitionv}{def:model}[Ambient model \(\mathcal M(\alpha,\beta)\)]
\(\mathcal M(\alpha,\beta)\) is the class of observed laws \(P\) satisfying the following conditions, for each \((\alpha,\beta)\in\mathcal D\), the prescribed domain of public smoothness exponents:
\begin{itemize}
\item \textbf{(Uniform design.)} With \(\lambda\) denoting the uniform probability law on the unit interval,
\[
P\circ X^{-1}=\lambda.
\]
See \cref{obj:ass:uniform-design}.
\item \textbf{(Homogeneous logit.)} The treated-arm conditional outcome risk \(\mu_1(P,x)\), logistic map \(\ell\), native prognosis logit \(\nu(P,x)\), and homogeneous conditional log-odds coefficient \(\theta(P)\) satisfy
\[
\mu_1(P,x)=\ell\{\nu(P,x)+\theta(P)\}
\qquad (x\in[0,1]).
\]
See \cref{obj:ass:homogeneous-logit}.
\item \textbf{(Effect envelope.)}
\[
|\theta(P)|\le \frac12.
\]
See \cref{obj:ass:effect-envelope}.
\item \textbf{(Propensity envelope.)} The native propensity-logit function \(g(P,\cdot)\) satisfies
\[
\|g(P,\cdot)\|_\infty\le 1.
\]
See \cref{obj:ass:propensity-envelope}.
\item \textbf{(Propensity smoothness.)}
\[
[g(P,\cdot)]_\alpha\le 2.
\]
See \cref{obj:ass:propensity-holder}.
\item \textbf{(Prognosis envelope.)}
\[
\|\nu(P,\cdot)\|_\infty\le 1.
\]
See \cref{obj:ass:prognosis-envelope}.
\item \textbf{(Prognosis smoothness.)}
\[
[\nu(P,\cdot)]_\beta\le 2.
\]
See \cref{obj:ass:prognosis-holder}.
\end{itemize}
The bracket notation denotes the Hölder seminorm at the indicated public exponent.
\end{definitionv}

\begin{assumptionv}{ass:evaluation-radius}[Effect evaluation radius]
The homogeneous conditional log-odds coefficient \(\theta(P)\) satisfies
\[
|\theta(P)|\le r.
\]
\end{assumptionv}

\begin{definitionv}{def:radius-model}[Effect-radius slice \(\mathcal M_r(\alpha,\beta)\)]
\[
\mathcal M_r(\alpha,\beta)
=
\{P\in\mathcal M(\alpha,\beta):|\theta(P)|\le r\}
\qquad
\bigl((\alpha,\beta)\in\mathcal D,\ r\in[0,1/2]\bigr).
\]
Here \(\mathcal D\) is the prescribed domain of public smoothness exponents and \(\theta(P)\) is the homogeneous conditional log-odds coefficient. This defines the slice of the ambient model in \cref{obj:def:model} satisfying \cref{obj:ass:evaluation-radius}.
\end{definitionv}

\begin{assumptionv}{ass:global-coverage}[Global interval coverage]
The interval procedure \(I(\mathcal O,U)\), evaluated on the tuple of observed records \(\mathcal O\) and an independent uniform randomization seed \(U\), satisfies
\[
\inf_{P\in\mathcal M(\alpha,\beta)}
(P^{\otimes n}\otimes\lambda)
\{\theta(P)\in I(\mathcal O,U)\}\ge 0.90,
\]
where \(\mathcal M(\alpha,\beta)\) is the ambient observed-law model, \(\lambda\) is the uniform probability law on the unit interval, and \(\theta(P)\) is the homogeneous conditional log-odds coefficient.
\end{assumptionv}

\begin{definitionv}{def:honest-procedures}[Honest interval procedures $\mathcal A_n(\alpha,\beta)$]
For every \(n\ge1\) and \((\alpha,\beta)\in\mathcal D\), define
\[
\mathcal A_n(\alpha,\beta)
=
\left\{
I:\Omega^n\times[0,1]\to\mathcal C:
\inf_{P\in\mathcal M(\alpha,\beta)}
(P^{\otimes n}\otimes\lambda)
\{\theta(P)\in I(\mathcal O,U)\}
\ge0.90
\right\},
\]
where the maps \(I\) are Borel and have measurable coverage events and interval length, as required by \cref{obj:ass:global-coverage}. Here \(\Omega\) is the measurable space of one observed record, \(\mathcal O\) is the tuple of original observed records, \(\mathcal C\) is the space of admissible nonempty connected Borel effect intervals, and \(\lambda\) is the uniform probability law on \([0,1]\). The variable \(U\) is an independent randomization seed, \(\mathcal D\) is the prescribed domain of public smoothness exponents, and \(\theta(P)\) is the homogeneous conditional log-odds coefficient.
\end{definitionv}

\begin{definitionv}{def:length-objective}[Minimax expected length $J_n(\alpha,\beta;r)$]
The minimax expected interval length on the effect-radius slice, subject to honesty over the ambient model, is
\[
J_n(\alpha,\beta;r)
=
\inf_{I\in\mathcal A_n(\alpha,\beta)}
\sup_{P\in\mathcal M_r(\alpha,\beta)}
E_{P^{\otimes n}\otimes\lambda}
\bigl|I(\mathcal O,U)\bigr|.
\]
Here \(\mathcal A_n(\alpha,\beta)\) is the class of Borel interval procedures with global \(90\%\) coverage, \(\mathcal M_r(\alpha,\beta)\) is the effect-radius slice, \(\mathcal O\) is the tuple of original observed records, and \(U\) is an independent seed with uniform law \(\lambda\) on \([0,1]\). The quantity \(\lvert I(\mathcal O,U)\rvert\) is the length of the returned interval.
\end{definitionv}

\begin{definitionv}{def:diagnostic-envelope}[Diagnostic envelope \(d_n(\alpha,\beta;r)\)]
The diagnostic envelope is
\[
d_n(\alpha,\beta;r)
=
n^{-\mathsf a(\alpha,\beta)}
+r\,n^{-\mathsf b(\beta)},
\]
where the numerator and denominator rate exponents are, respectively,
\[
\mathsf a(\alpha,\beta)
=
\min\left\{\frac12,\frac{2(\alpha+\beta)}{2\alpha+2\beta+1}\right\},
\qquad
\mathsf b(\beta)=\frac{4\beta}{4\beta+1}.
\]
\end{definitionv}

\begin{definitionv}{synth_3}[Cosine projection and its kernel]
Let \(\lambda\) be the uniform probability measure on \([0,1]\). The orthonormal cosine basis of \(L^2(\lambda)\) is
\[
\varphi_0(x)=1,\qquad
\varphi_j(x)=\sqrt{2}\cos(\pi jx)\quad(j\ge1).
\]
For each integer \(k\ge1\), define the rank-\(k\) orthogonal projection onto \(\operatorname{span}\{\varphi_0,\ldots,\varphi_{k-1}\}\) by
\[
\Pi_k f
=\sum_{j=0}^{k-1}\langle f,\varphi_j\rangle_{L^2(\lambda)}\varphi_j,
\qquad f\in L^2(\lambda),
\]
where \(\langle f,h\rangle_{L^2(\lambda)}=\int_0^1 f(z)h(z)\,d\lambda(z)\). Its projection kernel is
\[
H_k(x,z)=\sum_{j=0}^{k-1}\varphi_j(x)\varphi_j(z)
=1+2\sum_{j=1}^{k-1}\cos(\pi jx)\cos(\pi jz),
\qquad x,z\in[0,1].
\]
Thus the continuous representative of \(\Pi_k f\) satisfies
\[
(\Pi_k f)(x)=\int_0^1 H_k(x,z)f(z)\,d\lambda(z),
\qquad x\in[0,1].
\]
\end{definitionv}

\begin{definitionv}{def:projection-statistic}[Projection statistic $\mathbb U_{n,k}(W,V)$]
For \(n\ge2\), the ordered-pair projection statistic for two bounded record marks \(W\) and \(V\) is
\[
\mathbb U_{n,k}(W,V)
=
\frac{1}{n(n-1)}
\sum_{\substack{1\le i,j\le n\\ i\ne j}}
H_k(X_i,X_j)W(O_i)V(O_j).
\]
Here \(O_i\) is the \(i\)th observed record, \(X_i\) is its covariate, and \(H_k\) is the kernel of the projection at rank \(k\).
\end{definitionv}

\begin{definitionv}{def:dyadic-name-convention}[Dyadic naming policies \(\mathsf N,\mathsf N_0\)]
A dyadic name for a real number \(v\) is a sequence of nested closed intervals with dyadic rational endpoints that contain \(v\) and have width at most \(2^{-q}\) at each integer precision \(q\ge0\). The product coding of these sequences uses the discrete measurable structure on rational responses.

An admissible naming policy \(\mathsf N\) is a sequence of policies, one for each sample size \(n\), satisfying the following conditions.
\begin{itemize}
\item \textbf{(Public exponents.)} Valid names for \(\alpha,\beta\) are selected by Borel functions of the public pair \((\alpha,\beta)\) alone, before sampling.
\item \textbf{(Observed covariates.)} Valid names for the \(n\) observed covariates are selected by Borel functions of \((\alpha,\beta,\mathcal O)\), where \(\mathcal O\) is the tuple of original observed records.
\item \textbf{(Separate arithmetic input.)} The arithmetic engine is independently fixed, and name selection supplies the real-input representations used with that engine.
\end{itemize}
The interval procedure receives these selected names and the exact binary labels. Name selection is an input convention.

The dyadic-floor policy \(\mathsf N_0\) assigns every real input \(v\) the intervals
\[
\mathsf N_0(v)_q
=
\left[
2^{-q}\lfloor2^qv\rfloor,\;
2^{-q}\bigl(\lfloor2^qv\rfloor+1\bigr)
\right],
\qquad q\ge0.
\]
This is a valid name also at dyadic inputs and selects one concrete member of the policy-indexed interval family.

At the rank step, the queried exponent intervals are intersected with the known compact boxes
\[
\alpha\in[0,1],\qquad \beta\in[0,1/4].
\]
A syntactically malformed queried prefix or an empty intersection may yield the full effect interval. At the endpoint precision, the supplied exponent and observed-covariate intervals are passed raw to the endpoint expression tree, and the integer ranks selected at the rank step are retained. Outputs may vary across valid names.
\end{definitionv}

\begin{definitionv}{def:arithmetic-engine}[Arithmetic engine $\mathsf E$]
An arithmetic engine \(\mathsf E\) is a deterministic collection of maps on closed rational intervals, with integer precision input \(q\ge 0\) and tolerance
\[
\epsilon_q=2^{-q}.
\]
It supplies rational enclosures of \(\pi\), the exponential and cosine on every bounded rational interval, the logarithm on rational intervals \([a,b]\) with \(0<a\le b\), the square root on rational intervals \([a,b]\) with \(0\le a\le b\), and the power \(b^v\) for a positive integer base \(b\) and a rational exponent interval contained in \([-3,3]\). The symbols \(a,b,v,q\) are bound primitive arguments; in the power primitive, \(b\) denotes the integer base, supplied separately from the interval endpoints.

The engine also supplies exact rational range operations for addition, subtraction, multiplication, division on denominator intervals avoiding zero, minimum, maximum, intersection, and clipping. In particular, the minimum of two intervals has range
\[
\min\bigl([a,b],[c,d]\bigr)
=
[\min(a,c),\min(b,d)],
\]
and multiplication has range
\[
[a,b]\,[c,d]
=
[\min\{ac,ad,bc,bd\},\max\{ac,ad,bc,bd\}].
\]
Division uses the reciprocal range followed by multiplication, and clipping uses the exact minimum and maximum operations.

Admissibility requires the following primitive contracts.
\begin{itemize}
\item \textbf{(Enclosure.)} Every returned interval contains the real image of its input interval.
\item \textbf{(Monotone primitives.)} For exponential, logarithm, square root, and positive-integer-base power, each returned endpoint differs outward from the corresponding exact range endpoint by at most \(\epsilon_q\).
\item \textbf{(Constant enclosure.)} The returned enclosure of \(\pi\) lies in \([3,4]\), with outward excess at each endpoint at most \(\epsilon_q\).
\item \textbf{(Cosine.)} At the rational midpoint of the input interval, the engine encloses the cosine value with outward endpoint excess at most \(\epsilon_q\). It enlarges this enclosure by half the input width on each side and intersects the result with \([-1,1]\). The returned width is at most the input width plus \(2\epsilon_q\).
\item \textbf{(Primitive domains and signs.)} Exponential and power outputs have positive lower endpoints, and square-root outputs have nonnegative lower endpoints. Logarithm is evaluated on intervals with positive rational lower endpoints, and division is evaluated on denominator intervals avoiding zero.
\item \textbf{(Finite rational computation.)} Every map is defined by a finite rational recurrence whose iteration count is bounded by a specified public integer function of the rational range bounds, the integer base when present, and \(q\). The recurrence may use rational comparisons and integer arithmetic; its access to a real input is exclusively through the supplied rational interval.
\item \textbf{(Independent choice.)} The engine is fixed independently of all public exponents, sample sizes, observed records, laws, naming policies, evaluation radii, and unknown nuisance quantities.
\end{itemize}
Distinct admissible engines may return distinct rational enclosures. Admissibility is specified entirely by these primitive contracts, with rank conclusions, final-endpoint conclusions, final precision budgets, and statistical properties outside its defining requirements.
\end{definitionv}

\begin{algorithmv}{def:resolutions}[Public projection rank selection]
The projection ranks are the nonnegative integer ceilings of the upper endpoints produced by the following construction. For \(n\ge2\) and real public exponents \(\alpha,\beta\), define the public precision budget and target resolutions by
\[
q_{\mathrm{rank}}(n)=7+3\lceil\log_2 n\rceil,\qquad
R_C=n^{2/(2\alpha+2\beta+1)},\qquad
R_S=n^{2/(4\beta+1)}.
\]
For a positive integer base \(b\), the logarithmic ceiling is defined by integer comparisons:
\[
\lceil\log_2 b\rceil
=\min\{q\in\mathbb Z_{\ge0}:2^q\ge b\}.
\]

Fix a rational interval-arithmetic engine \(\mathsf E\) and a naming policy \(\mathsf N\). Query the public exponent names at precision \(q_{\mathrm{rank}}(n)\) and intersect them with \([0,1]\) and \([0,1/4]\), respectively. Using these intersections, form the ranges of
\[
\frac{2}{2\alpha+2\beta+1},
\qquad
\frac{2}{4\beta+1}
\]
by exact rational interval operations, then apply the power map of \(\mathsf E\) with integer base \(n\) at the same precision. If these operations succeed, denote the resulting boxes by \([a_C,b_C]\) and \([a_S,b_S]\), and define
\[
\begin{aligned}
k_C(n,\alpha,\beta;\mathsf E,\mathsf N)
&=\max\{\lceil b_C\rceil,0\},\\
k_S(n,\beta;\mathsf E,\mathsf N)
&=\max\{\lceil b_S\rceil,0\}.
\end{aligned}
\]
Here \(b_C,b_S\) denote resolution-box upper endpoints. If either an intersection or an interval division fails, define both ranks to be \(1\). Both ranks use the public inputs \((n,\alpha,\beta)\) through \(\mathsf N\), including in the abbreviated notation for \(k_S\), and are fixed before sampling.

For the admissible construction, take:
\begin{itemize}
\item \textbf{(Engine.)} \(\mathsf E\) satisfying \cref{obj:def:arithmetic-engine}.
\item \textbf{(Names.)} \(\mathsf N\) satisfying \cref{obj:def:dyadic-name-convention}, with public exponent names depending on \(n,\alpha,\beta\).
\item \textbf{(Exponent domain.)} \(0<\beta<1/4\) and \(\beta<\alpha\le1\).
\end{itemize}
For every \(n\ge2\), the construction succeeds and its boxes satisfy
\[
R_C\in[a_C,b_C],\qquad R_S\in[a_S,b_S],
\qquad b_C-a_C\le1,\qquad b_S-a_S\le1.
\]
Consequently, the selected ranks satisfy
\[
\begin{aligned}
R_C&\le k_C(n,\alpha,\beta;\mathsf E,\mathsf N)\le3R_C,\\
R_S&\le k_S(n,\beta;\mathsf E,\mathsf N)\le3R_S.
\end{aligned}
\]
\end{algorithmv}

\begin{definitionv}{def:coordinate-estimates}[Coordinate estimates $\widehat C_n$ and $\widehat S_n$]
The exact real sample coordinates are
\[
\widehat C_n
=\frac1n\sum_{i=1}^n A_iY_i-\mathbb U_{n,k_C}(A,Y),\qquad
\widehat S_n
=\mathbb U_{n,k_S}\bigl(A(1-Y),(1-A)Y\bigr),
\]
where, for \(n\ge2\), the retained ranks are
\[
k_C=k_C(n,\alpha,\beta;\mathsf E,\mathsf N),\qquad
k_S=k_S(n,\alpha,\beta;\mathsf E,\mathsf N).
\]
Both ranks are selected by the fixed-budget resolution construction with the same admissible engine \(\mathsf E\) and the public exponent approximations supplied by \(\mathsf N\). The upper construction uses this same engine and the supplied covariate approximations to enclose the finite sums at these retained ranks. The statistics \(\widehat C_n,\widehat S_n\) estimate the observable numerator and denominator coordinates \(C(P),S(P)\), respectively.
\end{definitionv}

\begin{definitionv}{def:causal-extension}[Full-data law $\widetilde P_{P,\Gamma}$]
The full-data law \(\widetilde P_{P,\Gamma}\) is defined by the successive draws
\[
X\sim\lambda,\qquad
(Y^0,Y^1)\mid X=x\sim\Gamma(P,x),\qquad
A\mid(X,Y^0,Y^1)\sim\operatorname{Bernoulli}\{e(P,X)\},
\]
and the observed outcome
\[
Y=(1-A)Y^0+AY^1.
\]
Here \(\lambda\) is the uniform probability law on \([0,1]\), \(Y^0\) and \(Y^1\) are the control and treated potential outcomes, \(\Gamma(P,x)\) is an admissible conditional coupling of these outcomes, and \(e(P,X)\) is the treatment probability conditional on the covariate under the observed law \(P\).
\end{definitionv}

\begin{definitionv}{synth_5}[Frontier comparison constants and sample-size threshold]
Fix positive finite absolute constants \(c,C\) satisfying the simultaneous comparisons in \cref{obj:thm:sharp-honest-frontier}. Define the selected lower and upper comparison functions and sample-size threshold on \(\mathcal D\) by
\[
c_\star(\alpha,\beta)=c,\qquad
C_\star(\alpha,\beta)=C,\qquad
n_0(\alpha,\beta)=1.
\]
Thus, for every admissible arithmetic engine \(\mathsf E\), admissible Borel naming policy \(\mathsf N\), \((\alpha,\beta)\in\mathcal D\), integer \(n\ge n_0(\alpha,\beta)\), and \(r\in[0,1/2]\),
\[
c_\star(\alpha,\beta)\rho_n(\alpha,\beta;r)
\le J_n(\alpha,\beta;r)
\le
\sup_{P\in\mathcal M_r(\alpha,\beta)}
E_{P^{\otimes n}}\!
\left|I^\star_{n,\mathsf E,\mathsf N}(\alpha,\beta)\right|
\le C_\star(\alpha,\beta)\rho_n(\alpha,\beta;r).
\]
The constants \(c,C\) are selected once, independently of the exponents, sample size, radius, engine, and naming policy.
\end{definitionv}

\begin{definitionv}{def:calibration-maps}[Calibration maps $\mathfrak c,\mathfrak B,\mathfrak q,\mathsf T,\mathfrak H,\mathfrak p$]
The covariance adjustment and its effect-normalized counterpart are
\[
\mathfrak c(t,\xi,\upsilon)
=\frac{2dv_*}{L_*+\sqrt{L_*^2-4d^2v_*}},\qquad
\mathfrak B(t,\xi,\upsilon)
=\frac{2D(t)v_*}{L_*+\sqrt{L_*^2-4d^2v_*}},
\]
where \(\xi,\upsilon\) are scalar margins near \(1/2\), and
\[
d=\exp(t)-1,\qquad
v_*=\xi(1-\xi)\upsilon(1-\upsilon),
\]
\[
L_*=1+d\{\xi(1-\upsilon)+(1-\xi)\upsilon\},\qquad
D(t)=\int_0^1\exp(st)\,ds.
\]
Here \(s\) is an integration variable. These formulas apply on the domain where the displayed square root and denominator are positive.

For two independent uniform signs \(\epsilon_0,\epsilon_1\), define
\[
Z_u=\epsilon_0\cos(\pi u/2)+\epsilon_1\sin(\pi u/2).
\]
Expectations below are over these signs. Let \(\varepsilon_{\mathrm{mix}}>0\) be the radius of the fixed signed open amplitude square on which the mixed calibration branch exists uniquely. On
\[
|\eta|<\varepsilon_{\mathrm{mix}},\qquad
|\zeta|<\varepsilon_{\mathrm{mix}},\qquad
u\in[0,1],
\]
define the mixed calibration root by
\[
\mathfrak q(\eta,\zeta,u)
=
\text{the unique }q\in(3/4,5/4)
\text{ satisfying }
16E\mathfrak B
\bigl(32\eta\zeta,1/2-\eta qZ_u,1/2+\zeta Z_u\bigr)-q=0.
\]
The smaller closed amplitude squares used for support and mixture bounds are contained in this domain.

For the separate fair-prior maps, define the baseline risk and local variance by
\[
p_*=2/5,\qquad v_*=p_*(1-p_*)=6/25.
\]
Here \(v_*\) denotes the fair-branch Bernoulli variance. Define
\[
\mathsf G(t,\xi)
=\frac{\exp(t)\xi}{1+\{\exp(t)-1\}\xi},\qquad
B_*=1+\{\exp(t)-1\}p_*,
\]
\[
\mathsf T(t,\delta)
=t+
\log\left(1-\frac{\{\exp(t)-1\}\delta^2}{p_*B_*}\right)
-\log\left(1+\frac{\{\exp(t)-1\}\delta^2}{(1-p_*)B_*}\right).
\]
On the fair core \(0\le t\le1/4\), for \(t\delta\ne0\), the normalized fair-matching residual is
\[
\mathfrak H(t,\delta,\xi,u)
=
\frac{
E\mathsf G(t,\xi+\delta Z_u)
-\mathsf G(\mathsf T(t,\delta),\xi)
}{t\delta^2}.
\]
The symbol \(\mathfrak H\) also denotes its removable smooth extension to the fixed open calibration neighborhood, including across \(t=0\), \(\delta=0\), and the spatial endpoints. The displayed quotient identity applies on the fair core.

Let \(\varepsilon_{\mathrm{fair}}>0\) be the signed fair-amplitude radius, and let \(b\in(0,1/10)\) be the fixed sufficiently small bracket half-width about \(p_*\). Here \(b\) denotes a bracket half-width. The fair parameter domain contains
\[
0\le t\le1/4,\qquad
|\delta|<\varepsilon_{\mathrm{fair}},\qquad
u\in[0,1],
\]
together with the fixed open neighborhood of these parameters, including signed parameters across \(t=0\) and the spatial endpoints. The denominators and logarithm arguments in the fair formulas are positive on this neighborhood. At \(t=0\) or \(\delta=0\), \(\mathfrak H\) is interpreted through its removable smooth extension.

On this neighborhood, define the fair calibration root by
\[
\mathfrak p(t,\delta,u)
=
\text{the unique }\xi\in(p_*-b,p_*+b)
\text{ satisfying }\mathfrak H(t,\delta,\xi,u)=0.
\]
Selection in the larger bracket \((3/10,1/2)\) gives this same root on the stated neighborhood. The smaller closed amplitude neighborhood used for support and mixture bounds is contained in this fair domain, uniformly over \(t\in[0,1/4]\) and \(u\in[0,1]\). The local branches are defined on the stated domains, where their existence, uniqueness, regularity, and matching properties hold.
\end{definitionv}

\begin{definitionv}{def:continuous-calibrated-priors}[Continuous calibrated priors \(P_{b,\sigma},P_{t,\sigma},Q_b,Q_t\)]
The continuous sign field on an equal partition is
\[
\mathsf Z_{k,\sigma}(x)
=
\sigma_j\cos(\pi u/2)+\sigma_{j+1}\sin(\pi u/2),
\qquad
u=\frac{x-j\Delta}{\Delta},
\qquad
\Delta=\frac1k,
\]
for an integer \(k\ge1\), independent uniform signs
\(\sigma=(\sigma_0,\ldots,\sigma_k)\), and
\(j\Delta\le x\le(j+1)\Delta\), \(0\le j<k\). Either adjoining cell gives the value at a shared endpoint.

For signed amplitudes \(\eta,\zeta\), define
\[
e_0(x)=\frac12+\eta\mathfrak q(\eta,\zeta,u)\mathsf Z_{k,\sigma}(x),
\qquad
e_1(x)=\frac12-\eta\mathfrak q(\eta,\zeta,u)\mathsf Z_{k,\sigma}(x),
\]
\[
m_\sigma(x)=\frac12+\zeta\mathsf Z_{k,\sigma}(x),
\qquad
c_0(x)=0,
\qquad
c_1(x)=\mathfrak c(32\eta\zeta,e_1(x),m_\sigma(x)).
\]
For the family index \(b\in\{0,1\}\), the law \(P_{b,\sigma}\) has uniform covariate law \(\lambda\) on the unit interval and conditional cell probabilities
\[
\begin{aligned}
p_{11}&=e_bm_\sigma+c_b,
&
p_{10}&=e_b(1-m_\sigma)-c_b,\\
p_{01}&=(1-e_b)m_\sigma-c_b,
&
p_{00}&=(1-e_b)(1-m_\sigma)+c_b.
\end{aligned}
\]
Here \(e_b,m_\sigma,c_b\) are evaluated at the current covariate; \(p_{11},p_{10},p_{01},p_{00}\) correspond to treated success, treated failure, control success, and control failure.

For \(0\le t\le1/4\), the separate fair-treatment laws \(P_{t,\sigma}\) and \(P_{*,t,\delta}\) have covariate law \(\lambda\), treatment probability \(1/2\), and control and treated outcome risks
\[
\begin{array}{c|cc}
&\mu_0(x)&\mu_1(x)\\ \hline
P_{t,\sigma}
&
\mathfrak p(t,\delta,u)+\delta\mathsf Z_{k,\sigma}(x)
&
\mathsf G\bigl(t,\mathfrak p(t,\delta,u)+\delta\mathsf Z_{k,\sigma}(x)\bigr)
\\
P_{*,t,\delta}
&
\mathfrak p(t,\delta,u)
&
\mathsf G\bigl(\mathsf T(t,\delta),\mathfrak p(t,\delta,u)\bigr).
\end{array}
\]
These formulas define probability laws on the calibration neighborhood where the conditional cells are positive.

The sample mixtures are
\[
Q_b
=
2^{-(k+1)}
\sum_{\sigma\in\{-1,1\}^{k+1}}
P_{b,\sigma}^{\otimes n},
\qquad
Q_t
=
2^{-(k+1)}
\sum_{\sigma\in\{-1,1\}^{k+1}}
P_{t,\sigma}^{\otimes n}.
\]
The signs are drawn once for the entire sample, and each mixture consists of exactly \(n\) original observed records.
\end{definitionv}

\begin{definitionv}{def:frontier-handle}[Rate profile $\rho$ and attaining interval family $I^\star$]
The rate profile and its attaining interval family are defined, for every admissible engine \(\mathsf E\) and admissible naming policy \(\mathsf N\), by
\[
\rho_n(\alpha,\beta;r)=d_n(\alpha,\beta;r),
\qquad
I^\star_{n,\mathsf E,\mathsf N}
=
I^{\mathrm{up}}_{n,\mathsf E,\mathsf N}.
\]
Here \(d_n(\alpha,\beta;r)\) is the diagnostic expected-length envelope from \cref{obj:def:diagnostic-envelope}, and \(I^{\mathrm{up}}_{\mathsf E,\mathsf N}\) is the procedure sequence from \cref{obj:def:upper-handle} for the fixed engine and policy. The frontier pair consists of this rate profile and procedure sequence:
\[
\bigl(\rho,I^{\mathrm{up}}_{\mathsf E,\mathsf N}\bigr).
\]
The concrete attaining sequence is
\[
I_n^\star
=
I^{\mathrm{up}}_{n,\mathsf E_0,\mathsf N_0},
\]
where \(\mathsf E_0\) is the explicitly constructed rational interval-arithmetic engine supplied by \cref{obj:lem:concrete-arithmetic-engine}, and \(\mathsf N_0\) is the dyadic-floor policy for selecting numerical representations in \cref{obj:def:dyadic-name-convention}.

Each interval is the literal output of the same fixed-budget expression tree, evaluated by its supplied fixed engine on its supplied valid numerical representations. Each choice of engine and representations determines its own rational output, subject to the endpoint enclosure bounds in \cref{obj:lem:fixed-budget-realization}. The construction uses the observable numerator and denominator coordinates \(C(P),S(P)\) from \cref{obj:lem:observable-odds}, separate projection resolutions from \cref{obj:def:resolutions}, and joint inversion from \cref{obj:def:upper-handle}.

The represented inputs are \((n,\alpha,\beta,\mathcal O)\), with numerical representations selected by \(\mathsf N\) and an independently fixed engine \(\mathsf E\) as in \cref{obj:def:arithmetic-engine}. The radius \(r\) indexes evaluation of the expected-length profile.

The converse uses the continuous calibrated priors in \cref{obj:def:continuous-calibrated-priors}, with admissible native propensity and prognosis perturbations, one common scalar effect per draw, analytic singleton matching, and factorization across disjoint latent-sign components conditional on the original covariates. It applies independently of the engine and naming policy.

The model and honesty conditions are specified in \cref{obj:synth_1,obj:synth_2,obj:synth_3,obj:ass:uniform-design,obj:ass:homogeneous-logit,obj:ass:effect-envelope,obj:ass:propensity-envelope,obj:ass:prognosis-envelope,obj:ass:propensity-holder,obj:ass:prognosis-holder,obj:ass:evaluation-radius,obj:ass:global-coverage,obj:def:logistic-class,obj:def:model,obj:def:radius-model,obj:def:honest-procedures,obj:def:length-objective}. The associated constructions are given in \cref{obj:def:causal-extension,obj:def:projection-statistic,obj:def:coordinate-estimates,obj:def:calibration-maps}. Their bounds, calibrations, and sharp comparisons are stated in \cref{obj:lem:projection-control,obj:thm:finite-honest-upper,obj:lem:parametric-null-floor,obj:thm:sharp-honest-frontier,obj:lem:scalar-implicit-function,obj:lem:exact-calibrations,obj:lem:calibrated-support,obj:lem:original-record-mixtures,obj:thm:full-honest-frontier-solution,obj:thm:full-honest-frontier-answer}.
\end{definitionv}

\begin{theoremv}{thm:full-honest-frontier-answer}[The full honest frontier]
There exist constants \(c,C\in\mathbb R\), with \(c>0\) and \(C>0\), satisfying the following guarantees simultaneously. The public exponent domain and the sharp expected-length profile are
\[
\mathcal D=\{(\alpha,\beta)\in\mathbb R^2:0<\beta<1/4,\ \beta<\alpha\le1\},
\]
\[
\rho_n(\alpha,\beta;r)=d_n(\alpha,\beta;r)
=n^{-\min\{1/2,\,2(\alpha+\beta)/(2\alpha+2\beta+1)\}}
+r\,n^{-4\beta/(4\beta+1)}.
\]

\begin{itemize}
\item \textbf{(Concrete attainment.)}
For every \((\alpha,\beta)\in\mathcal D\) and every integer \(n\ge1\), the concrete interval
\[
I_n^\star(\alpha,\beta)
=I^{\mathrm{up}}_{n,\mathsf E_0,\mathsf N_0}(\alpha,\beta),
\]
constructed with the explicit rational arithmetic engine \(\mathsf E_0\) and the dyadic-floor naming policy \(\mathsf N_0\), belongs to the globally honest class \(\mathcal A_n(\alpha,\beta)\) of \cref{obj:def:honest-procedures}. For every sample and randomization seed, its output is a closed interval \([l,u]\) with
\[
l,u\in\mathbb Q,\qquad -1/2\le l\le u\le1/2.
\]
For every \(r\in[0,1/2]\), the minimax expected length \(J_n(\alpha,\beta;r)\) of \cref{obj:def:length-objective} satisfies
\[
c\,\rho_n(\alpha,\beta;r)
\le J_n(\alpha,\beta;r)
\le
\sup_{P\in\mathcal M_r(\alpha,\beta)}
E_{P^{\otimes n}\otimes\lambda}
\bigl|I_n^\star(\alpha,\beta)(\mathcal O,U)\bigr|
\le C\,\rho_n(\alpha,\beta;r),
\]
where \(\mathcal M_r(\alpha,\beta)\) is the prescribed radius slice and \(\lambda\) is the uniform law of the independent randomization seed.

\item \textbf{(Universal construction.)}
For every admissible rational interval-arithmetic engine \(\mathsf E\), every admissible Borel naming policy \(\mathsf N\), both satisfying the quantitative contracts of \cref{obj:def:arithmetic-engine,obj:def:dyadic-name-convention}, every \((\alpha,\beta)\in\mathcal D\), and every integer \(n\ge1\), the fixed-budget output
\(I^{\mathrm{up}}_{n,\mathsf E,\mathsf N}(\alpha,\beta)\)
is Borel, globally \(90\%\)-honest, and radius independent. On every valid named input, the construction terminates within its prescribed fixed budgets and returns a nonempty closed interval with ordered rational endpoints in \([-1/2,1/2]\), using the observed data and public inputs. For every \(r\in[0,1/2]\),
\[
c\,\rho_n(\alpha,\beta;r)
\le J_n(\alpha,\beta;r)
\le
\sup_{P\in\mathcal M_r(\alpha,\beta)}
E_{P^{\otimes n}\otimes\lambda}
\bigl|I^{\mathrm{up}}_{n,\mathsf E,\mathsf N}(\alpha,\beta)(\mathcal O,U)\bigr|
\le C\,\rho_n(\alpha,\beta;r).
\]
The same constants \(c,C\) apply to all these engines, policies, exponents, sample sizes, and radii.

\item \textbf{(Compact uniformity.)}
Taking the comparison constants \(c_\star=c\), \(C_\star=C\), and the sample-size threshold \(n_0=1\), every nonempty compact set \(K\subseteq\mathcal D\) satisfies
\[
0<\inf_{(\alpha,\beta)\in K}c_\star,\qquad
\sup_{(\alpha,\beta)\in K}C_\star<\infty,\qquad
\sup_{(\alpha,\beta)\in K}n_0=1.
\]

\item \textbf{(Elbow identities.)}
Define the numerator and denominator exponents by
\[
\mathsf a(\alpha,\beta)
=\min\{1/2,\,2(\alpha+\beta)/(2\alpha+2\beta+1)\},
\qquad
\mathsf b(\beta)=\frac{4\beta}{4\beta+1}.
\]
For every \((\alpha,\beta)\in\mathcal D\),
\[
\mathsf a(\alpha,\beta)-\mathsf b(\beta)>0,
\]
and
\[
\mathsf a(\alpha,\beta)-\mathsf b(\beta)
=
\begin{cases}
\displaystyle
\frac{2(\alpha-\beta)}
{(2\alpha+2\beta+1)(4\beta+1)},
&\alpha+\beta\le1/2,\\[6pt]
\displaystyle
\frac{1-4\beta}{2(4\beta+1)},
&\alpha+\beta\ge1/2.
\end{cases}
\]
At \(\alpha+\beta=1/2\), the two displayed expressions are equal.

\item \textbf{(Universal converse.)}
For every \((\alpha,\beta)\in\mathcal D\), every \(r\in[0,1/2]\), every integer \(n\ge1\), and every procedure \(I\in\mathcal A_n(\alpha,\beta)\), including procedures using independent randomization,
\[
c\,\rho_n(\alpha,\beta;r)
\le
\sup_{P\in\mathcal M_r(\alpha,\beta)}
E_{P^{\otimes n}\otimes\lambda}
\bigl|I(\mathcal O,U)\bigr|.
\]
\end{itemize}
\end{theoremv}

\begin{algorithmv}{def:upper-handle}[Upper interval $I^{\mathrm{up}}_{n,\mathsf E,\mathsf N}(\alpha,\beta)$]
The inputs are the sample size \(n\ge1\), public smoothness exponents \(\alpha,\beta\), and observed sample \(\mathcal O\); the tuning choices are an independently fixed admissible rational interval-arithmetic engine \(\mathsf E\) and a Borel policy \(\mathsf N\) selecting valid numerical representations of the exponents and observed covariates.
\begin{enumerate}
\item For \(n=1\), return
\[
I^{\mathrm{up}}_{1,\mathsf E,\mathsf N}(\alpha,\beta)
=[-1/2,1/2].
\]
For \(n\ge2\), proceed through the remaining steps.

\item Select the numerator and denominator projection ranks using \(\mathsf E\), the public exponent representations selected by \(\mathsf N\), and the fixed public precision budget \(q_{\mathrm{rank}}(n)\):
\[
k_C=k_C(n,\alpha,\beta;\mathsf E,\mathsf N),
\qquad
k_S=k_S(n,\alpha,\beta;\mathsf E,\mathsf N).
\]
Use these retained integers in the coordinate estimates \(\widehat C_n\), the numerator estimate, and \(\widehat S_n\), the denominator estimate, from \cref{obj:def:coordinate-estimates}, with ranks defined in \cref{obj:def:resolutions}.

\item Define the smoothness sum, variance envelope, and coordinate error radii by
\[
s=\alpha+\beta,\qquad
V_n(k)=\frac{10}{n}+\frac{4k}{n(n-1)},
\]
\[
b_C=8k_C^{-s}+\sqrt{20V_n(k_C)},
\qquad
b_S=25k_S^{-2\beta}+\sqrt{20V_n(k_S)}.
\]
Here \(b_C\) and \(b_S\) denote error radii, distinct from the resolution-enclosure endpoints in \cref{obj:def:resolutions}.

\item Define scalar clipping, the fixed denominator lower bound from \cref{obj:lem:observable-odds}, and the confidence-rectangle endpoints by
\[
\operatorname{clip}_{[l,u]}(v)=\max\{l,\min\{u,v\}\},
\qquad s_0=\frac{1}{512},
\]
\[
c_\pm=\widehat C_n\pm b_C,
\qquad
s_\pm=\operatorname{clip}_{[s_0,1]}(\widehat S_n\pm b_S).
\]
Define the clipped logarithmic inversion and its ideal interval endpoints by
\[
F(c,d)=
\log\!\left(
\operatorname{clip}_{[\exp(-1/2),\exp(1/2)]}(1+c/d)
\right),
\]
\[
L=\min\{F(c_-,s_-),F(c_-,s_+)\},
\qquad
R=\max\{F(c_+,s_-),F(c_+,s_+)\}.
\]
Here \(R\) denotes the ideal upper endpoint, distinct from the generic rank target in \cref{obj:def:resolutions}. These exact-real expressions specify the quantities enclosed in the next step.

\item Define the largest retained rank and the endpoint precision budget by
\[
K=\max\{k_C,k_S\},
\qquad
q_{\mathrm{end}}(n,k_C,k_S)
=44+\lceil\log_2(n+1)\rceil
+4\lceil\log_2(K+1)\rceil.
\]
At this single precision, query the supplied numerical representations of both exponents and all observed covariates. Evaluate the displayed finite expressions with the same engine \(\mathsf E\), its primitive maps at precision \(q_{\mathrm{end}}\), and exact rational range operations in the displayed expression tree, retaining \(k_C,k_S\) from the rank step. Obtain rational endpoint enclosures satisfying
\[
L_-\le L\le L_+,
\qquad
R_-\le R\le R_+.
\]
For every valid input with \(0<\beta<1/4\), \(\beta<\alpha\le1\), and \(n\ge2\), \cref{obj:lem:fixed-budget-realization} gives
\[
L_+-L_-\le\frac1n,
\qquad
R_+-R_-\le\frac1n.
\]

\item Output the closed rational interval
\[
I^{\mathrm{up}}_{n,\mathsf E,\mathsf N}(\alpha,\beta)
=
\left[
\max\{-1/2,L_-\},
\min\{1/2,R_+\}
\right].
\]
For every valid input with \(0<\beta<1/4\), \(\beta<\alpha\le1\), and \(n\ge2\), the output is nonempty, contains the ideal interval, and satisfies
\[
[L,R]\subseteq
I^{\mathrm{up}}_{n,\mathsf E,\mathsf N}(\alpha,\beta),
\]
\[
0\le
\min\{1/2,R_+\}-\max\{-1/2,L_-\}-(R-L)
\le\frac2n.
\]
The computation terminates after the fixed rank and endpoint precision budgets.
\end{enumerate}
For fixed \(\mathsf E,\mathsf N\), this construction defines a sample map with inputs \((n,\alpha,\beta,\mathcal O)\). Its rational endpoints may depend on the admissible engine and valid numerical representations supplied.

The model and honesty conditions are specified in \cref{obj:synth_1,obj:synth_2,obj:synth_3,obj:ass:uniform-design,obj:ass:homogeneous-logit,obj:ass:effect-envelope,obj:ass:propensity-envelope,obj:ass:prognosis-envelope,obj:ass:propensity-holder,obj:ass:prognosis-holder,obj:ass:evaluation-radius,obj:ass:global-coverage,obj:def:logistic-class,obj:def:model,obj:def:radius-model,obj:def:honest-procedures,obj:def:length-objective}. The associated constructions are given in \cref{obj:def:causal-extension,obj:def:projection-statistic,obj:def:diagnostic-envelope,obj:def:frontier-handle,obj:def:calibration-maps,obj:def:continuous-calibrated-priors,obj:def:dyadic-name-convention,obj:def:arithmetic-engine}. Their bounds and realization properties are stated in \cref{obj:lem:projection-control,obj:thm:finite-honest-upper,obj:lem:parametric-null-floor,obj:thm:sharp-honest-frontier,obj:lem:scalar-implicit-function,obj:lem:exact-calibrations,obj:lem:calibrated-support,obj:lem:original-record-mixtures,obj:thm:full-honest-frontier-solution,obj:lem:concrete-arithmetic-engine,obj:thm:full-honest-frontier-answer}.
\end{algorithmv}

\begin{theoremv}{thm:finite-honest-upper}[Uniform honesty and length bound]
There exists an absolute constant \(K_0>0\) such that the following guarantees hold for every choice satisfying:
\begin{itemize}
\item \textbf{(Arithmetic engine.)} \(\mathsf E\) is an admissible, independently fixed deterministic rational interval-arithmetic engine satisfying the quantitative contracts of \cref{obj:def:arithmetic-engine}.
\item \textbf{(Naming policy.)} \(\mathsf N\) is an admissible Borel policy satisfying \cref{obj:def:dyadic-name-convention}, selecting valid names for the public exponents and observed covariates, with public exponent names fixed independently of the sample.
\item \textbf{(Public exponents.)} The real exponents \(\alpha,\beta\) satisfy
\[
0<\beta<\frac14,\qquad \beta<\alpha\le1.
\]
\item \textbf{(Sample size.)} The integer sample size \(n\) satisfies \(n\ge1\).
\end{itemize}

Let \(I^{\mathrm{up}}_{n,\mathsf E,\mathsf N}(\alpha,\beta)\) denote the measurable interval procedure given by the literal rational output of the upper construction, whose inputs are the engine, naming policy, sample size, public exponents, and observed records. This procedure satisfies
\[
I^{\mathrm{up}}_{n,\mathsf E,\mathsf N}(\alpha,\beta)
\in\mathcal A_n(\alpha,\beta),
\]
where \(\mathcal A_n(\alpha,\beta)\) is the class of globally honest procedures in \cref{obj:def:honest-procedures}.

For every \(r\in[0,1/2]\) and every observed law \(P\in\mathcal M_r(\alpha,\beta)\), the effect-radius model of \cref{obj:def:radius-model},
\[
E_{P^{\otimes n}\otimes\lambda}
\left[
\left|I^{\mathrm{up}}_{n,\mathsf E,\mathsf N}(\alpha,\beta)(\mathcal O,U)\right|
\right]
\le K_0\,d_n(\alpha,\beta;r).
\]
Here \(\mathcal O\) is the sample of original observed records, \(U\) is the independent randomization seed with uniform law \(\lambda\) on the unit interval, and \(d_n(\alpha,\beta;r)\) is the diagnostic envelope of \cref{obj:def:diagnostic-envelope}. The same constant \(K_0\) applies to all the engines, naming policies, exponents, sample sizes, radii, and laws quantified above.

For every sample and seed, there exist rational endpoints \(l,u\) such that the output is the closed interval \([l,u]\) and
\[
-\frac12\le l\le u\le\frac12.
\]
When \(n=1\), every sample and seed yields the full interval \([-1/2,1/2]\).
\end{theoremv}

\begin{lemmav}{lem:observable-odds}[Observable odds and causal interpretation]
Let \(0<\beta<1/4\) and \(\beta<\alpha\leq1\), and let \(P\in\mathcal M(\alpha,\beta)\), the model class in \cref{obj:def:model}. Write \(\lambda\) for the uniform probability law on \([0,1]\), and define the observable coordinates
\[
C(P)
=
E_P[AY]
-
\int_{[0,1]} e(P,x)\,E_P[Y\mid X=x]\,\lambda(dx),
\qquad
S(P)
=
\int_{[0,1]}
P(A=1,Y=0\mid X=x)\,
P(A=0,Y=1\mid X=x)\,\lambda(dx),
\]
where \(e(P,x)\) is the conditional treatment probability. With
\[
\theta(P)=\operatorname{logit}\mu_1(P,0)-\nu(P,0),
\qquad s_0=\frac{1}{512},
\]
where \(\mu_1(P,0)\) is the treated-arm risk at zero and \(\nu(P,0)\) is the control-arm prognosis logit at zero, one has
\[
C(P)=\{\exp(\theta(P))-1\}S(P),
\qquad
S(P)\geq s_0,
\qquad
\theta(P)=\log\{1+C(P)/S(P)\}.
\]

For every Markov kernel \(\Gamma\) from \([0,1]\) to \(\{0,1\}^2\) whose first and second margins at every \(x\) are Bernoulli laws with means
\[
P(Y=1\mid A=0,X=x)
\quad\text{and}\quad
P(Y=1\mid A=1,X=x),
\]
respectively, the full-data law \(\widetilde P_{P,\Gamma}\) of \cref{obj:def:causal-extension} has the following properties:
\begin{itemize}
\item \textbf{(Observed law and consistency.)}
The image of \(\widetilde P_{P,\Gamma}\) under
\[
((x,(y^0,y^1)),a)\longmapsto
\bigl(x,a,(1-a)y^0+ay^1\bigr)
\]
is \(P\). For every full-data record,
\[
Y=(1-A)Y^0+AY^1.
\]

\item \textbf{(Construction and conditional exchangeability.)}
The law \(\widetilde P_{P,\Gamma}\) is the successive kernel composition obtained by drawing \(X\) from \(\lambda\), drawing \((Y^0,Y^1)\) from \(\Gamma(P,X)\), and drawing \(A\) from the observed-law treatment kernel with Bernoulli mean \(e(P,X)\). Under its full conditional law at every \(x\in[0,1]\), the pair \((Y^0,Y^1)\) and \(A\) are independent.

\item \textbf{(Conditional causal log-odds.)}
For every \(x\in[0,1]\),
\[
\operatorname{logit}
\left(\int_{\{0,1\}^2} y^1\,\Gamma(P,x)(d(y^0,y^1))\right)
-
\operatorname{logit}
\left(\int_{\{0,1\}^2} y^0\,\Gamma(P,x)(d(y^0,y^1))\right)
=
\theta(P).
\]

\item \textbf{(Full conditional disintegration.)}
Composing \(\lambda\) with the full conditional laws of \(((Y^0,Y^1),A)\) at \(x\), then taking the image under the canonical reassociation
\[
\bigl(x,((y^0,y^1),a)\bigr)
\longmapsto
\bigl((x,(y^0,y^1)),a\bigr),
\]
gives exactly \(\widetilde P_{P,\Gamma}\).
\end{itemize}
\end{lemmav}

\begin{lemmav}{lem:projection-control}[Projection approximation and moments]
Let \(P\) be an observed law with continuous, strictly interior conditional four-cell probabilities, satisfying \cref{obj:ass:uniform-design}. Let \(k\ge1\) be an integer, and let \(\Pi_k\) denote the rank-\(k\) cosine projection used in \cref{obj:def:projection-statistic}.

For every \(0<\gamma\le1\) and every continuous function \(f:[0,1]\to\mathbb R\),
\[
\|f-\Pi_k f\|_{L^2(\lambda)}
\le 5[f]_\gamma k^{-\gamma},
\qquad
[f]_\gamma
=
\sup_{\substack{x,z\in[0,1]\\x\ne z}}
\frac{|f(x)-f(z)|}{|x-z|^\gamma}.
\]
Here \(\lambda\) is the uniform probability law on \([0,1]\), and the norm and Hölder seminorm are interpreted in the extended nonnegative reals, allowing \([f]_\gamma=+\infty\).

For every pair of measurable record marks \(W,V\) taking values in \([0,1]\), and every integer \(n\ge2\), define their conditional means using the four-cell conditional law under \(P\):
\[
w(x)=E_P[W\mid X=x],
\qquad
v(x)=E_P[V\mid X=x].
\]
For the ordered-pair statistic \(\mathbb U_{n,k}(W,V)\) of \cref{obj:def:projection-statistic}, evaluated on \(n\) independent original records with law \(P\),
\[
E_P\mathbb U_{n,k}(W,V)
=
\langle\Pi_k w,\Pi_k v\rangle_{L^2(\lambda)},
\]
\[
\langle w,v\rangle_{L^2(\lambda)}
-
E_P\mathbb U_{n,k}(W,V)
=
\langle w-\Pi_k w,v-\Pi_k v\rangle_{L^2(\lambda)},
\]
and
\[
\operatorname{Var}_P\{\mathbb U_{n,k}(W,V)\}
\le
\frac4n+\frac{2k}{n(n-1)}.
\]
\end{lemmav}

\begin{lemmav}{lem:concrete-arithmetic-engine}[Admissibility of the concrete engine]
The explicit rational arithmetic engine \(\mathsf E_0\), consisting of the Machin enclosure for \(\pi\), the cosine interval routine, and the endpoint extensions of the scalar exponential, logarithm, square-root, and integer-base power routines, satisfies all primitive admissibility contracts of \cref{obj:def:arithmetic-engine}.
\end{lemmav}

\begin{lemmav}{lem:fixed-budget-realization}[Fixed-budget realization and measurability]
The dyadic-floor naming policy \(\mathsf N_0\) is admissible in the sense of \cref{obj:def:dyadic-name-convention}. Moreover, suppose that:
\begin{itemize}
\item \textbf{(Engine and names.)} \(\mathsf E\) is an admissible rational interval-arithmetic engine in the sense of \cref{obj:def:arithmetic-engine}, and \(\mathsf N\) is an admissible Borel naming policy in the sense of \cref{obj:def:dyadic-name-convention}.
\item \textbf{(Sample size.)} \(n\) is an integer with \(n\ge2\).
\item \textbf{(Public exponents.)} The real smoothness exponents satisfy
\[
0<\beta<\frac14,\qquad \beta<\alpha\le1.
\]
\end{itemize}
For every record tuple \(\mathcal O\in([0,1]\times\{0,1\}\times\{0,1\})^n\), let \(k_C,k_S\) be the ranks selected according to \cref{obj:def:resolutions}, with target resolutions
\[
R_C=n^{2/(2\alpha+2\beta+1)},\qquad
R_S=n^{2/(4\beta+1)}.
\]
The prescribed public budgets are
\[
q_{\mathrm{rank}}(n)=7+3\lceil\log_2 n\rceil,\qquad
q_{\mathrm{end}}(n,k_C,k_S)
=44+\lceil\log_2(n+1)\rceil
 +4\lceil\log_2(K+1)\rceil,\qquad
K=\max\{k_C,k_S\}.
\]
At the rank budget, both rational target enclosures \([a_C,b_C]\) and \([a_S,b_S]\) are successfully returned and satisfy
\[
b_C-a_C\le1,\qquad b_S-a_S\le1,
\]
\[
R_C\le k_C<R_C+2\le3R_C,\qquad
R_S\le k_S<R_S+2\le3R_S.
\]
Here \(b_C,b_S\) denote the upper endpoints of the resolution enclosures.

For the endpoint construction, use \(b_C,b_S\) instead to denote the coordinate error radii
\[
b_C=8k_C^{-(\alpha+\beta)}+\sqrt{20V_n(k_C)},\qquad
b_S=25k_S^{-2\beta}+\sqrt{20V_n(k_S)},
\]
where \(V_n(k)\) is the public variance envelope at rank \(k\). Let \(\widehat C_n,\widehat S_n\) be the numerator and denominator coordinate estimates, and set
\[
c_\pm=\widehat C_n\pm b_C,\qquad
s_\pm=\operatorname{clip}_{[s_0,1]}(\widehat S_n\pm b_S),
\]
where \(s_0\) is the denominator floor. With \(F(c,d)\) denoting the clipped logarithmic inversion of the coordinate ratio, the ideal lower and upper endpoints are
\[
L=\min\{F(c_-,s_-),F(c_-,s_+)\},\qquad
R=\max\{F(c_+,s_-),F(c_+,s_+)\}.
\]
At the endpoint budget, rational enclosures \([L_-,L_+]\) and \([R_-,R_+]\) are successfully returned with
\[
L\in[L_-,L_+],\qquad R\in[R_-,R_+],\qquad
L_+-L_-\le\frac1n,\qquad R_+-R_-\le\frac1n.
\]
The literal output satisfies
\[
[L,R]\subseteq
I^{\mathrm{up}}_{n,\mathsf E,\mathsf N}(\alpha,\beta)(\mathcal O),
\qquad
\operatorname{length}\!\left(
I^{\mathrm{up}}_{n,\mathsf E,\mathsf N}(\alpha,\beta)(\mathcal O)
\right)
\le R-L+\frac2n.
\]
It is a nonempty closed connected interval with rational endpoints in \([-1/2,1/2]\), and its output map, as a function of the record tuple, is Borel measurable.
\end{lemmav}

\begin{lemmav}{lem:scalar-implicit-function}[Scalar implicit function theorem]
Let \(d\) be a nonnegative integer, let \(F:\mathbb R^d\times\mathbb R\to\mathbb R\), and let \(D\subseteq\mathbb R^d\times\mathbb R\). Fix \(\mathbf b_0\in\mathbb R^d\) and \(\chi_0\in\mathbb R\). Suppose that:
\begin{itemize}
\item \textbf{(Open domain.)} \(D\) is open and \((\mathbf b_0,\chi_0)\in D\).
\item \textbf{(Smoothness.)} \(F\) is continuously differentiable on \(D\).
\item \textbf{(Zero.)} \(F(\mathbf b_0,\chi_0)=0\).
\item \textbf{(Nonzero partial derivative.)} \(\partial_\chi F(\mathbf b_0,\chi_0)\ne0\).
\end{itemize}
Then there exist open sets \(U\subseteq\mathbb R^d\) and \(V\subseteq\mathbb R\), and a function \(q:\mathbb R^d\to\mathbb R\), such that
\[
\mathbf b_0\in U,\qquad \chi_0\in V,\qquad U\times V\subseteq D.
\]
The function \(q\) is continuously differentiable on \(U\), satisfies \(q(\mathbf b_0)=\chi_0\), and has \(q(\mathbf b)\in V\) for every \(\mathbf b\in U\). Moreover, for every \(\mathbf b\in U\) and every \(\chi\in V\),
\[
F(\mathbf b,\chi)=0
\quad\Longleftrightarrow\quad
\chi=q(\mathbf b).
\]
\end{lemmav}

\begin{lemmav}{lem:exact-calibrations}[Exact smooth calibration branches]
There exist constants \(\varepsilon>0\) and \(M>0\) with the following properties. Expectations below are over two independent fair signs, and
\[
Z_u=s_1\cos(\pi u/2)+s_2\sin(\pi u/2),
\qquad
E F(s_1,s_2)=\frac14\sum_{s_1,s_2\in\{-1,1\}}F(s_1,s_2).
\]
The covariance adjustment \(\mathfrak c\) and its normalized extension \(\mathfrak B\) are given by
\[
\begin{gathered}
d=e^t-1,\qquad
v_*=\xi(1-\xi)\upsilon(1-\upsilon),\qquad
L_*=1+d\{\xi(1-\upsilon)+(1-\xi)\upsilon\},\\
D(t)=\int_0^1 e^{rt}\,dr,\\
\mathfrak c(t,\xi,\upsilon)
=\frac{2d v_*}{L_*+\sqrt{L_*^2-4d^2v_*}},
\qquad
\mathfrak B(t,\xi,\upsilon)
=\frac{2D(t)v_*}{L_*+\sqrt{L_*^2-4d^2v_*}}.
\end{gathered}
\]
The mixed root \(\mathfrak q(\eta,\zeta,u)\) is selected from the roots in \((3/4,5/4)\) of
\[
16E\mathfrak B
 \bigl(32\eta\zeta,\tfrac12-\eta qZ_u,\tfrac12+\zeta Z_u\bigr)-q=0,
\]
with value \(1\) when that bracket contains no root. The fair root
\(\mathfrak p(t,\delta,u)\) is selected from the roots in \((3/10,1/2)\) of
\(\mathfrak H(t,\delta,p,u)=0\), with value \(2/5\) when that bracket contains no root. Here \(\mathfrak H\) is the smoothly extended normalized fair-matching residual, whose defining identity away from the axes is displayed below.

\begin{itemize}
\item \textbf{(Positive odds tables.)}
For every \(t,\xi,\upsilon\in\mathbb R\) satisfying
\[
|t|\le\varepsilon,\qquad
|\xi-\tfrac12|\le\varepsilon,\qquad
|\upsilon-\tfrac12|\le\varepsilon,
\]
put \(c=\mathfrak c(t,\xi,\upsilon)\). The four cells
\[
\begin{aligned}
p_{11}&=\xi\upsilon+c,&
p_{10}&=\xi(1-\upsilon)-c,\\
p_{01}&=(1-\xi)\upsilon-c,&
p_{00}&=(1-\xi)(1-\upsilon)+c
\end{aligned}
\]
are strictly positive and satisfy
\[
p_{11}+p_{10}=\xi,\qquad
p_{11}+p_{01}=\upsilon,\qquad
p_{11}+p_{10}+p_{01}+p_{00}=1,\qquad
\frac{p_{11}p_{00}}{p_{10}p_{01}}=e^t.
\]
Moreover,
\[
\mathfrak B(t,\xi,\upsilon)=\frac{\mathfrak c(t,\xi,\upsilon)}{t}
\quad(t\ne0),\qquad
\mathfrak B(0,\xi,\upsilon)=\xi(1-\xi)\upsilon(1-\upsilon).
\]

\item \textbf{(Signed mixed calibration.)}
For every \(|\eta|\le\varepsilon\), \(|\zeta|\le\varepsilon\), and \(u\in[0,1]\), put
\(q=\mathfrak q(\eta,\zeta,u)\). Then
\[
\frac34<q<\frac54,\qquad
16E\mathfrak B
 \bigl(32\eta\zeta,\tfrac12-\eta qZ_u,\tfrac12+\zeta Z_u\bigr)-q=0,
\]
and
\[
\begin{gathered}
\mathfrak q(\eta,\zeta,0)=\mathfrak q(\eta,\zeta,1),\\
E\mathfrak c
 \bigl(32\eta\zeta,\tfrac12-\eta qZ_u,\tfrac12+\zeta Z_u\bigr)
 =2\eta\zeta q,\\
\mathfrak q(-\eta,\zeta,u)=q,\qquad
\mathfrak q(\eta,-\zeta,u)=q,\\
\mathfrak q(0,\zeta,u)=1-4\zeta^2,\qquad
\mathfrak q(\eta,0,u)+4\eta^2\mathfrak q(\eta,0,u)^2=1.
\end{gathered}
\]

\item \textbf{(Fair calibration and separation.)}
For every \(t\in[0,1/4]\), \(|\delta|\le\varepsilon\), and \(u\in[0,1]\), put
\(p=\mathfrak p(t,\delta,u)\). Then
\[
\begin{gathered}
\mathfrak p(t,\delta,0)=\mathfrak p(t,\delta,1)=\frac25,\qquad
\mathfrak p(t,0,u)=\mathfrak p(0,\delta,u)=\frac25,\\
\mathfrak p(t,-\delta,u)=p,\qquad
\mathfrak H(t,\delta,p,u)=0,\\
E\mathsf G(t,p+\delta Z_u)=\mathsf G(\mathsf T(t,\delta),p).
\end{gathered}
\]
For \(t\delta\ne0\), the residual at the selected root agrees with the unnormalized fair-matching numerator divided by \(t\delta^2\). Furthermore,
\[
|p-\tfrac25|+|\partial_u\mathfrak p(t,\delta,u)|\le M\delta^2,
\]
and
\[
3t\delta^2\le t-\mathsf T(t,\delta)\le6t\delta^2,\qquad
\frac t2\le\mathsf T(t,\delta)\le t.
\]

\item \textbf{(Smooth neighborhoods and derivative bounds.)}
Both roots and every fully substituted local four-cell formula of the mixed and fair families are \(C^\infty\) on their respective closed parameter regions
\[
\begin{gathered}
|\eta|\le\varepsilon,\quad |\zeta|\le\varepsilon,\quad u\in[0,1],\\
t\in[0,1/4],\quad |\delta|\le\varepsilon,\quad u\in[0,1].
\end{gathered}
\]
For every integer \(m\ge0\), there is a constant \(B>0\) such that the norm of the \(m\)th total derivative of each root and each fully substituted local four-cell formula is at most \(B\) throughout its respective closed parameter region, uniformly over all family indices, auxiliary signs, and binary cell indices.

There also exist signed amplitude radii
\(\varepsilon_{\mathrm{mix}}>\varepsilon\) and
\(\varepsilon_{\mathrm{fair}}>\varepsilon\), a fixed fair-root bracket half-width
\(0<b<1/10\), and open sets \(U,V\subset\mathbb R^3\) containing, respectively,
\[
\begin{gathered}
\{(\eta,\zeta,u):|\eta|\le\varepsilon_{\mathrm{mix}},
                       \ |\zeta|\le\varepsilon_{\mathrm{mix}},\ u\in[0,1]\},\\
\{(t,\delta,u):t\in[0,1/4],\
                  |\delta|\le\varepsilon_{\mathrm{fair}},\ u\in[0,1]\}.
\end{gathered}
\]
Throughout \(U\), the selected mixed branch is the unique zero in its fixed bracket \((3/4,5/4)\); throughout \(V\), the selected fair branch is the unique zero in the fixed centered bracket \((2/5-b,2/5+b)\). Both roots and every fully substituted local four-cell formula are \(C^\infty\) on the corresponding open set.

Finally, there is an open set \(W\subset\mathbb R^4\) containing
\[
\{(t,\delta,\xi,u):t\in[0,1/4],\
 |\delta|\le\varepsilon_{\mathrm{fair}},\
 \xi\in(2/5-b,2/5+b),\ u\in[0,1]\}
\]
on which \(\mathfrak H\) is \(C^\infty\). Throughout this displayed core, for every margin \(\xi\) in the centered bracket and whenever \(t\delta\ne0\),
\[
\mathfrak H(t,\delta,\xi,u)
=\frac{E\mathsf G(t,\xi+\delta Z_u)-\mathsf G(\mathsf T(t,\delta),\xi)}
       {t\delta^2}.
\]
\end{itemize}
\end{lemmav}

\begin{lemmav}{lem:calibrated-support}[Uniform calibrated support and smoothness]
There exist jointly chosen absolute constants \(\varepsilon\), the calibration radius, and \(M\), the derivative bound, satisfying
\[
0<\varepsilon\le \frac{1}{100},
\qquad M\ge 1,
\]
with the following properties for the calibrated families of \cref{obj:def:continuous-calibrated-priors}. Throughout, \(k\ge1\) is an integer and \(\sigma\in\{-1,1\}^{k+1}\) is a sign vector.

\begin{itemize}
\item \textbf{(Mixed membership and effects.)}
For every \((\alpha,\beta)\in\mathcal D\), the prescribed domain of public smoothness exponents, set
\[
\eta=\varepsilon k^{-\alpha},
\qquad
\zeta=\varepsilon k^{-\beta}.
\]
The calibrated mixed laws \(P_{0,\sigma}\) and \(P_{1,\sigma}\) at these amplitudes belong to \(\mathcal M(\alpha,\beta)\) as defined in \cref{obj:def:model}, and their homogeneous conditional log-odds coefficients satisfy
\[
\theta(P_{0,\sigma})=0,
\qquad
\theta(P_{1,\sigma})=32\eta\zeta.
\]
For each treatment label, outcome label, and covariate \(x\), the sums over all sign vectors of the corresponding conditional cell probabilities under the two mixed laws are equal.

\item \textbf{(Fair membership and effects.)}
For every \((\alpha,\beta)\in\mathcal D\), set
\[
\delta=\varepsilon k^{-\beta}.
\]
For every \(t\in[0,1/4]\), the fair random law \(P_{t,\sigma}\) and the deterministic fair comparator \(P_{*,t,\delta}\) belong to \(\mathcal M(\alpha,\beta)\), and
\[
\theta(P_{t,\sigma})=t,
\qquad
\theta(P_{*,t,\delta})=\mathsf T(t,\delta),
\]
where \(\mathsf T(t,\delta)\) is the calibrated comparator effect defined in \cref{obj:def:calibration-maps}. For each treatment label, outcome label, and covariate \(x\), the uniform average over the \(2^{k+1}\) sign vectors of the conditional cell probability under \(P_{t,\sigma}\) equals the corresponding comparator cell probability.

\item \textbf{(Signed validity and derivative bounds.)}
For all signed amplitudes satisfying
\[
|\eta|\le\varepsilon,
\qquad
|\zeta|\le\varepsilon,
\qquad
|\delta|\le\varepsilon,
\]
the conditional four-cell tables of both mixed branches are valid probability tables, as are those of both fair branches for every \(t\in[0,1/4]\). For every sign vector, branch, treatment label, outcome label, and covariate \(x\), their one-record densities relative to uniform four-label measure lie in \([1/2,2]\). The mixed densities, viewed as functions of \((\eta,\zeta)\), and the fair densities, viewed as functions of \(\delta\) with \(t\) fixed, are twice continuously differentiable at every point of the corresponding closed amplitude box. The norms of their first and second total amplitude derivatives are at most \(M\), uniformly over all these indices and amplitudes. These bounds include the mixed second derivative in \(\eta\) and \(\zeta\) and the deterministic comparator's derivatives in \(\delta\).

\item \textbf{(Common smooth neighborhoods.)}
There exists an open subset of \(\mathbb R^2\) containing
\[
\{(\eta,\zeta):|\eta|\le\varepsilon,\ |\zeta|\le\varepsilon\}
\]
on which every mixed density is infinitely continuously differentiable in its amplitude arguments. There also exists an open subset of \(\mathbb R\) containing
\[
[-\varepsilon,\varepsilon]
\]
on which every fair density is infinitely continuously differentiable in \(\delta\). Each neighborhood is common to every integer \(k\ge1\), sign vector, branch, treatment label, outcome label, and covariate \(x\); the fair neighborhood is also common to every \(t\in[0,1/4]\).

\item \textbf{(Signed singleton matching.)}
For every \(|\eta|\le\varepsilon\) and \(|\zeta|\le\varepsilon\), and each treatment label, outcome label, and covariate \(x\), the sums over all sign vectors of the corresponding conditional cell probabilities under the two mixed laws are equal. For every \(|\delta|\le\varepsilon\) and \(t\in[0,1/4]\), the uniform average over all \(2^{k+1}\) sign vectors of each conditional cell probability under the fair random law equals the corresponding cell probability under the deterministic fair comparator.
\end{itemize}
\end{lemmav}

\begin{definitionv}{synth_6}[Squared Hellinger distance for original-record sample laws]
For probability laws \(Q,Q'\) on the original-record sample space \(\Omega^n\), define their unhalved squared Hellinger distance by
\[
\mathcal H^2(Q,Q')
=
\int_{\Omega^n}
\left(
\sqrt{\frac{dQ}{dR}}
-
\sqrt{\frac{dQ'}{dR}}
\right)^2\,dR,
\]
where \(R\) is any common dominating measure. The value is independent of the choice of \(R\) and lies in \([0,2]\). For the continuous calibrated mixtures, take \(R\) to be the original-record sample reference measure. This normalization assigns coefficient one to the displayed integral.
\end{definitionv}

\begin{lemmav}{lem:original-record-mixtures}[Original-record mixture bounds]
Let \(\varepsilon\) and \(M\) be the radius and envelope constant selected from the fixed joint calibrated-support certificate, with \(0<\varepsilon\le 1/100\) and \(M\ge 1\). Define the positive constants
\[
\kappa=\frac{1}{96\exp(1)},
\qquad
C_{\mathrm{mix}}=96\exp(2)\,8^2M^4.
\]
For sample laws \(Q,Q'\), use the unhalved squared Hellinger distance
\[
\mathcal H^2(Q,Q')
=
\int
\left(
\sqrt{\frac{dQ}{dR}}-
\sqrt{\frac{dQ'}{dR}}
\right)^2\,dR,
\]
where \(R\) is the original-record sample reference measure.

For every choice satisfying
\begin{itemize}
\item \textbf{(Sample and partition sizes.)} \(n,k\) are integers with \(n\ge 1\) and \(k\ge 1\).
\item \textbf{(Occupancy.)} \(n/k\le\kappa\).
\item \textbf{(Signed amplitudes.)} \(\eta,\zeta,\delta\in\mathbb R\) satisfy
\[
|\eta|\le\varepsilon,\qquad
|\zeta|\le\varepsilon,\qquad
|\delta|\le\varepsilon,
\]
\end{itemize}
let \(Q_b\), for \(b\in\{0,1\}\), and \(Q_t\) be the uniform finite mixtures of the calibrated mixed and fair original-record sample laws, respectively:
\[
Q_b=\mathbb E_\sigma\!\left[P_{b,\sigma}^{\otimes n}\right],
\qquad
Q_t=\mathbb E_\sigma\!\left[P_{t,\sigma}^{\otimes n}\right],
\]
where the expectation averages uniformly over the finite sign vectors. Let \(P_{*,t,\delta}\) be the deterministic fair comparator obtained from the fair construction at partition size \(k\), center \(t\), and amplitude \(\delta\), using the constant numerical sign vector with every entry equal to \( -1 \). Then
\[
\mathcal H^2(Q_0,Q_1)
\le
C_{\mathrm{mix}}\frac{n^2}{k}\eta^2\zeta^2,
\]
and, for every \(t\in[0,1/4]\),
\[
\mathcal H^2\!\left(Q_t,P_{*,t,\delta}^{\otimes n}\right)
\le
C_{\mathrm{mix}}\frac{n^2}{k}\delta^4.
\]
\end{lemmav}

\begin{lemmav}{lem:parametric-null-floor}[Parametric lower bound on length]
Let \(\alpha,\beta,r\in\mathbb R\) and \(n\in\mathbb N\) satisfy:
\begin{itemize}
\item \textbf{(Exponent domain.)} \((\alpha,\beta)\in\mathcal D\), the public exponent domain given by
\[
\mathcal D=\{(\alpha,\beta)\in\mathbb R^2:0<\beta<1/4,\ \beta<\alpha\le1\}.
\]
\item \textbf{(Sample size.)} \(n\ge1\).
\item \textbf{(Radius.)} \(0\le r\le1/2\).
\end{itemize}
Then the minimax expected interval length \(J_n(\alpha,\beta;r)\) defined in \cref{obj:def:length-objective} satisfies
\[
J_n(\alpha,\beta;r)\ge \frac{3}{16}n^{-1/2}.
\]
\end{lemmav}

\begin{theoremv}{thm:sharp-honest-frontier}[Sharp honest interval frontier]
There exist positive finite absolute constants \(c,C\) with the following guarantees. Write
\[
\mathcal D=\{(\alpha,\beta)\in\mathbb R^2:0<\beta<1/4,\ \beta<\alpha\le1\},
\]
and define the numerator and denominator exponents and the expected-length profile by
\[
\mathsf a(\alpha,\beta)
=\min\left\{\frac12,\frac{2(\alpha+\beta)}{2\alpha+2\beta+1}\right\},
\qquad
\mathsf b(\beta)=\frac{4\beta}{4\beta+1},
\qquad
\rho_n(\alpha,\beta;r)
=n^{-\mathsf a(\alpha,\beta)}+r n^{-\mathsf b(\beta)}.
\]

For every admissible rational interval-arithmetic engine \(\mathsf E\) and every admissible Borel policy \(\mathsf N\) satisfying the quantitative contracts of \cref{obj:def:arithmetic-engine,obj:def:dyadic-name-convention} and selecting valid names for the public exponents and observed covariates, consider parameters satisfying:
\begin{itemize}
\item \textbf{(Exponent domain.)} \((\alpha,\beta)\in\mathcal D\).
\item \textbf{(Sample size.)} \(n\in\mathbb N\) and \(n\ge1\).
\item \textbf{(Radius.)} \(0\le r\le1/2\).
\end{itemize}
Let \(I^\star_{n,\mathsf E,\mathsf N}(\alpha,\beta)
=I^{\mathrm{up}}_{n,\mathsf E,\mathsf N}(\alpha,\beta)\) denote the interval returned by the fixed-budget upper construction. With \(\mathcal A_n(\alpha,\beta)\) as in \cref{obj:def:honest-procedures}, define the minimax expected length by
\[
J_n(\alpha,\beta;r)
=
\inf_{I\in\mathcal A_n(\alpha,\beta)}
\sup_{P\in\mathcal M_r(\alpha,\beta)}
\mathbb E_{P^{\otimes n}\otimes\lambda}
\bigl|I(\mathcal O,U)\bigr|,
\]
where \(\mathcal M_r(\alpha,\beta)\) is the prescribed effect-radius slice, \(\mathcal O\) is the sample of original records, and \(U\) is an independent seed with uniform law \(\lambda\) on the unit interval. Then
\[
c\rho_n(\alpha,\beta;r)
\le J_n(\alpha,\beta;r)
\le
\sup_{P\in\mathcal M_r(\alpha,\beta)}
\mathbb E_{P^{\otimes n}}
\bigl|I^\star_{n,\mathsf E,\mathsf N}(\alpha,\beta)\bigr|
\le C\rho_n(\alpha,\beta;r).
\]
The same constants apply to every admissible engine and naming policy.

For every such engine and policy, every \((\alpha,\beta)\in\mathcal D\), and every \(n\ge1\), the actual name-dependent output is a deterministic, radius-independent Borel interval with global \(90\%\) coverage over \(\mathcal M(\alpha,\beta)\), the prescribed ambient model. It is a nonempty closed connected interval contained in \([-1/2,1/2]\), with rational endpoints, returned after the two prescribed finite precision budgets and finite rational arithmetic. Its inputs consist of the fixed engine, the public exponents and their valid public names, and the \(n\) original records and their valid covariate names. Fixing the explicit admissible engine \(\mathsf E_0\) and dyadic-floor naming policy \(\mathsf N_0\) gives the concrete attaining interval
\[
I_n^\star
=I^{\mathrm{up}}_{n,\mathsf E_0,\mathsf N_0},
\qquad
c_\star=c,\qquad C_\star=C,\qquad n_0=1,
\]
where \(c_\star,C_\star\) are the comparison constants and \(n_0\) is the sample-size threshold.

Moreover, for every \((\alpha,\beta)\in\mathcal D\), every \(0\le r\le1/2\), every \(n\ge1\), and every \(I\in\mathcal A_n(\alpha,\beta)\), including procedures using independent randomization,
\[
c\rho_n(\alpha,\beta;r)
\le
\sup_{P\in\mathcal M_r(\alpha,\beta)}
\mathbb E_{P^{\otimes n}\otimes\lambda}
\bigl|I(\mathcal O,U)\bigr|.
\]

For every \((\alpha,\beta)\in\mathcal D\), the exponent difference is positive and satisfies
\[
\mathsf a(\alpha,\beta)-\mathsf b(\beta)>0,
\qquad
\mathsf a(\alpha,\beta)-\mathsf b(\beta)
=
\begin{cases}
\displaystyle
\frac{2(\alpha-\beta)}
{(2\alpha+2\beta+1)(4\beta+1)},
&\alpha+\beta\le1/2,\\[6pt]
\displaystyle
\frac{1-4\beta}{2(4\beta+1)},
&\alpha+\beta\ge1/2.
\end{cases}
\]
At \(\alpha+\beta=1/2\), the two displayed expressions agree.

Finally, for every nonempty compact set \(K\subseteq\mathcal D\), the same constants and threshold satisfy the compact-uniform conditions
\[
0<\inf_{(\alpha,\beta)\in K}c,
\qquad
\sup_{(\alpha,\beta)\in K}C<\infty,
\qquad
\sup_{(\alpha,\beta)\in K}1=1.
\]
\end{theoremv}

\begin{theoremv}{thm:full-honest-frontier-solution}[Full honest frontier solution]
There exist absolute constants \(c,C\in\mathbb R\), with \(0<c\) and \(0<C\), for which the following assertions hold. Write
\[
\mathcal D=\{(\alpha,\beta)\in\mathbb R^2:0<\beta<1/4,\ \beta<\alpha\le1\},
\qquad
\mathsf a(\alpha,\beta)=
\min\left\{\frac12,\frac{2(\alpha+\beta)}{2\alpha+2\beta+1}\right\},
\qquad
\mathsf b(\beta)=\frac{4\beta}{4\beta+1},
\]
and define the expected-length profile by
\[
\rho_n(\alpha,\beta;r)
=d_n(\alpha,\beta;r)
=n^{-\mathsf a(\alpha,\beta)}
+r n^{-\mathsf b(\beta)}.
\]

For every admissible rational interval-arithmetic engine \(\mathsf E\), every admissible Borel naming policy \(\mathsf N\), both satisfying the quantitative contracts of \cref{obj:def:arithmetic-engine,obj:def:dyadic-name-convention}, every \((\alpha,\beta)\in\mathcal D\), and every integer \(n\ge1\), take
\[
I^\star_{n,\mathsf E,\mathsf N}(\alpha,\beta)
=I^{\mathrm{up}}_{n,\mathsf E,\mathsf N}(\alpha,\beta),
\]
the literal output of the fixed-budget represented upper construction. This procedure belongs to the globally \(90\%\)-honest Borel class \(\mathcal A_n(\alpha,\beta)\) of \cref{obj:def:honest-procedures}. It receives neither the radius nor unknown nuisance functions and terminates on every valid named input within the prescribed rank and endpoint budgets of \cref{obj:def:resolutions,obj:def:upper-handle,obj:lem:fixed-budget-realization}. For every sample and randomization seed its output is a closed interval \([l,u]\) with
\[
l,u\in\mathbb Q,
\qquad -\frac12\le l\le u\le\frac12.
\]
For every \(r\in[0,1/2]\), the same constants satisfy
\[
c\rho_n(\alpha,\beta;r)
\le J_n(\alpha,\beta;r)
\le
\sup_{P\in\mathcal M_r(\alpha,\beta)}
E_P\!\left[
\left|I^\star_{n,\mathsf E,\mathsf N}(\alpha,\beta)(\mathcal O,U)\right|
\right]
\le C\rho_n(\alpha,\beta;r).
\]
Here \(J_n(\alpha,\beta;r)\) is the minimax expected length of \cref{obj:def:length-objective}, and the expectation uses \(n\) original records drawn independently from \(P\), together with the independent randomization seed \(U\).

In particular, the explicit concrete engine \(\mathsf E_0\) and the dyadic-floor naming policy \(\mathsf N_0\) give the sequence
\[
I_n^\star(\alpha,\beta)
=I^{\mathrm{up}}_{n,\mathsf E_0,\mathsf N_0}(\alpha,\beta).
\]
For every \((\alpha,\beta)\in\mathcal D\) and \(n\ge1\), this sequence is globally honest, has the rational closed ordered output described above at every input, and satisfies the same displayed comparison for every \(r\in[0,1/2]\).

Moreover, for every \((\alpha,\beta)\in\mathcal D\), every \(r\in[0,1/2]\), every integer \(n\ge1\), and every \(I\in\mathcal A_n(\alpha,\beta)\), including procedures using independent randomization,
\[
c\rho_n(\alpha,\beta;r)
\le
\sup_{P\in\mathcal M_r(\alpha,\beta)}
E_P\!\left[|I(\mathcal O,U)|\right].
\]

The exponent difference is positive throughout \(\mathcal D\) and obeys
\[
\mathsf a(\alpha,\beta)-\mathsf b(\beta)=
\begin{cases}
\displaystyle
\frac{2(\alpha-\beta)}
{(2\alpha+2\beta+1)(4\beta+1)},
&\alpha+\beta\le1/2,\\[6pt]
\displaystyle
\frac{1-4\beta}{2(4\beta+1)},
&\alpha+\beta\ge1/2.
\end{cases}
\]
At \(\alpha+\beta=1/2\), the two displayed expressions agree.

Finally, select the lower and upper comparison constants and sample-size threshold as
\[
c_\star(\alpha,\beta)=c,\qquad
C_\star(\alpha,\beta)=C,\qquad
n_0(\alpha,\beta)=1.
\]
For every nonempty compact \(\mathcal K\subset\mathcal D\), these choices satisfy
\[
0<
\inf_{(\alpha,\beta)\in\mathcal K}c_\star(\alpha,\beta),
\qquad
\sup_{(\alpha,\beta)\in\mathcal K}C_\star(\alpha,\beta)<\infty,
\qquad
\sup_{(\alpha,\beta)\in\mathcal K}n_0(\alpha,\beta)=1.
\]
\end{theoremv}
